\documentclass[journal, a4paper]{IEEEtran}
\usepackage[utf8]{vietnam}
\usepackage{cite}
\usepackage{amsmath,amssymb,amsfonts}
\usepackage{graphicx}
\usepackage{textcomp,nicefrac}
\usepackage{booktabs}
\usepackage{tabularx}
\usepackage{multirow}
\usepackage{float}

\def\BibTeX{{\rm B\kern-.05em{\sc i\kern-.025em b}\kern-.08em
T\kern-.1667em\lower.7ex\hbox{E}\kern-.125emX}}
\markboth{STUDENT SCIENTIFIC RESEARCH COMMUNICATION, Volume II, 2026}
{Nguyen \MakeLowercase{\textit{et al.}}: Design and Implementation of a Personalized Healthcare System}

\begin{document}

\title{HealthCare Now: Design and Implementation of a Personalized Healthcare System Powered by Artificial Intelligence}

\author{Nguyen Cong Danh and Tran Quoc Bao, \IEEEmembership{Faculty of Software Engineering, IUH}
\thanks{This manuscript was submitted for review in May 2026. This research was conducted by students at the Faculty of Information Technology, Industrial University of Ho Chi Minh City (IUH), Vietnam.}
\thanks{N. C. Danh and T. Q. Bao are with the Industrial University of Ho Chi Minh City (e-mail: 22635131.danh@student.iuh.edu.vn).}}

\maketitle

\begin{abstract}
In the era of digital health, the shift from reactive medicine to proactive, personalized care is an objective necessity, driven by the rise of chronic diseases and the demand for higher quality of life. However, current mobile health (mHealth) applications face significant limitations, including scattered Patient-Generated Health Data (PGHD), lack of deep personalization, and the underutilization of Artificial Intelligence (AI). This paper presents the exhaustive design, development, and evaluation of "HealthCare Now," a comprehensive health management system. The study leverages a robust Microservices Architecture orchestrating four core domains via an Nginx API Gateway, utilizing Node.js (BFF), Spring Boot (Core \& IoT), and FastAPI (AI). It implements event-driven communication via RabbitMQ using the Saga Pattern to ensure distributed medical data consistency, alongside Polyglot Persistence using MongoDB and Redis. To address data fragmentation, the system strategically incorporates standard health storage models. The core AI module is deeply integrated with biomedical theories, utilizing the Mifflin-St Jeor metabolic equations for nutritional baselines. Analytically, the platform employs a hybrid Machine Learning approach: XGBoost and Random Forest for dynamic physiological predictions; deep learning models (LSTM Autoencoders and Isolation Forests) for time-series anomaly detection; CNN-LSTM pipelines for OCR digitization; and Retrieval-Augmented Generation (RAG) powered by Google Gemini Large Language Models (LLM) to completely eliminate hallucinations in medical consultations. The evaluation results demonstrate that the FastAPI-driven ASGI architecture vastly outperforms WSGI frameworks in serving ML models, achieving high predictive clinical accuracy, robust scalability, real-time processing capabilities, and strict compliance with data privacy frameworks.

Trong kỷ nguyên y tế số, sự chuyển dịch từ y học phản ứng sang chăm sóc sức khỏe chủ động, cá nhân hóa là một tất yếu khách quan, được thúc đẩy bởi sự gia tăng của các bệnh mãn tính. Tuy nhiên, các ứng dụng mHealth hiện tại đối mặt với những hạn chế lớn, bao gồm dữ liệu phân tán (PGHD), thiếu tính cá nhân hóa và chưa tận dụng triệt để AI. Bài báo này trình bày toàn diện việc thiết kế, phát triển và đánh giá hệ thống "HealthCare Now". Nghiên cứu sử dụng Kiến trúc Microservices mạnh mẽ điều phối bốn miền cốt lõi thông qua Nginx, kết hợp Node.js, Spring Boot và FastAPI. Hệ thống áp dụng giao tiếp hướng sự kiện qua RabbitMQ và lưu trữ Polyglot Persistence với MongoDB và Redis. Về mặt phân tích, nền tảng áp dụng Học máy kết hợp: XGBoost cho dự đoán sinh lý động; LSTM Autoencoders để phát hiện bất thường chuỗi thời gian; CNN-LSTM cho số hóa OCR; và RAG kết hợp LLM để loại bỏ ảo giác trong tư vấn y tế. Kết quả đánh giá cho thấy kiến trúc ASGI dựa trên FastAPI vượt trội hơn hẳn các framework WSGI, đạt độ chính xác lâm sàng cao, khả năng mở rộng mạnh mẽ và khả năng xử lý thời gian thực xuất sắc.
\end{abstract}

\begin{IEEEkeywords}
Artificial Intelligence, FastAPI, IoT, LLM, mHealth, Microservices, RAG, Spring Boot, Time-Series Anomaly Detection, XGBoost.
\end{IEEEkeywords}

\section{Introduction}
\label{sec:introduction}

\IEEEPARstart{T}{he} explosion of information technology, the Internet of Things (IoT), and artificial intelligence (AI) over the past decade has fundamentally reshaped how humans manage their health. The transition from reactive medicine—where treatment only occurs after a pathology has manifested—to proactive and personalized medicine is an objective and inevitable trend.

\subsection{Background and Problem Statement}
The core driver for this shift stems from immense epidemiological pressures on a global scale. According to statistical reports by the World Health Organization (WHO) for the 2024-2025 period, non-communicable diseases (NCDs) such as cardiovascular diseases, cancer, chronic respiratory diseases, and diabetes are the leading causes of death, claiming approximately 41 million lives annually, equivalent to 74\% of all deaths worldwide.

In Vietnam, the epidemiological landscape poses even more severe challenges. Recent statistics indicate that the mortality rate from non-communicable diseases in Vietnam has reached an alarming threshold, accounting for 77\% to over 80\% of all annual deaths, exceeding the global average. A significant proportion of these are premature deaths, primarily resulting from sedentary lifestyles, unbalanced nutrition, and the lack of a proactive health monitoring system. The rise in chronic conditions necessitates lifelong treatment, creating a massive financial burden that directly threatens the sustainability of the national Health Insurance fund.

In response to this public health crisis, macro-level policies in Vietnam are strongly promoting digital transformation. A prime example is Ho Chi Minh City's strategic vision up to 2030, which emphasizes building a generation of "Digital Citizens" by leveraging the Government's Project 06. The healthcare system aims to fully integrate Electronic Health Records (EHR) into the national identification application (VNeID), while interconnecting tens of millions of specialized data records to personalize care services and establish early warning models for every citizen.

\subsection{Bottlenecks of Contemporary mHealth Systems}
Within this modern digital ecosystem, Patient-Generated Health Data (PGHD) continuously collected from smart wearables acts as the core lifeblood, providing multi-dimensional biometric streams such as heart rate, step count, sleep quality, and heart rate variability. However, simply amassing a vast volume of raw data does not automatically translate into clinical benefits. Through an extensive survey of existing health platforms, several critical bottlenecks have been identified:

\begin{enumerate}
    \item \textbf{Superficial Tracking Capabilities:} Most current applications act merely as data loggers. Users can view their daily step counts, but the systems fail to analyze deep historical data to provide a comprehensive evaluation. They cannot effectively identify whether a user's health trajectory is improving or degrading over time.
    \item \textbf{Limited Personalization:} Despite AI integration in various fields, its application in personal healthcare remains generalized. Current systems often provide static, generic advice rather than tailoring recommendations based on an individual's unique biological makeup. Applying a "one-size-fits-all" evaluation model severely limits optimal clinical outcomes.
    \item \textbf{Underutilized Artificial Intelligence:} Most apps only display statistics or trigger simple threshold-based alerts. AI's true capability—such as analyzing complex time-series data, detecting subtle physiological anomalies, predicting health trends, and generating personalized nutritional regimens—remains largely unexploited.
    \item \textbf{Scattered and Unsynchronized Data:} Health data is heavily fragmented. Movement data is stored in Google Fit, wearable metrics in Garmin, dietary logs in standalone apps, and clinical records exist as physical papers. This fragmentation prevents the creation of a holistic health profile necessary for training personalized AI models.
\end{enumerate}

\subsection{Research Purpose and Objectives}
To address these critical shortcomings, this thesis presents the exhaustive design, development, and evaluation of "HealthCare Now". The overall objective is to architect a smart, personalized health management system that solves data fragmentation and leverages a microservices-driven AI engine to provide proactive, evidence-based healthcare interventions. The detailed objectives include:
\begin{itemize}
    \item To construct a highly scalable backend using a 4-domain Microservices Architecture that mitigates the bottlenecks of monolithic designs.
    \item To embed biomedical algorithms (e.g., Mifflin-St Jeor) to dynamically calculate user metabolic rates and macronutrient distribution.
    \item To train and deploy Machine Learning models (XGBoost, Random Forest) for predicting dynamic physiological targets.
    \item To implement a Medical Chatbot utilizing Retrieval-Augmented Generation (RAG) combined with Google Gemini LLMs to eradicate AI "hallucinations."
    \item To apply deep learning (CNN-LSTM-BioBERT) to automate the digitization of paper-based medical records via Optical Character Recognition (OCR).
\end{itemize}

\subsection{Clinical Scope and Target Population}
Technically, this research heavily focuses on the backend architecture, data engineering, and the training, optimization, and deployment of AI models. Clinically, the proposed "HealthCare Now" system specifically targets the prevention and management of Non-Communicable Diseases (NCDs) associated with Metabolic Syndrome and Cardiovascular Health. This primarily encompasses cardiovascular risk monitoring (e.g., detecting heart rate anomalies such as tachycardia via LSTM Autoencoders and predicting cardiac responses via XGBoost) and metabolic disorders (e.g., preventing obesity and managing Type 2 Diabetes through dynamic caloric tracking and the Mifflin-St Jeor metabolic equations). The mobile frontend developed via React Native serves to validate these interventions and handle user interactions, rather than acting as the core subject of standalone diagnostic claims.

The current implementation of "HealthCare Now" primarily focuses on preventative healthcare, cardiovascular monitoring, metabolic health assessment, and lifestyle-related physiological analysis rather than direct disease diagnosis. Specifically, the system is optimized for:
\begin{itemize}
    \item cardiovascular risk monitoring through heart rate and activity analysis,
    \item calorie expenditure and metabolic tracking,
    \item sleep and recovery assessment,
    \item anomaly detection in wearable physiological signals,
    \item and personalized nutritional recommendations based on real-time biometric data.
\end{itemize}
The proposed AI models are intended to function as decision-support and early-warning tools for wellness management. They are not designed to replace professional medical diagnosis or certified clinical procedures. Although the architecture is extensible to broader medical domains, the current study concentrates on chronic lifestyle-related conditions and physiological monitoring scenarios commonly associated with non-communicable diseases (NCDs), including cardiovascular disorders, obesity, metabolic syndrome, and sedentary lifestyle risks.

\section{Literature Review and Theoretical Foundations}
\label{sec:theory}

Building a comprehensive health monitoring system cannot rely solely on software algorithms but must be anchored in dynamic biomedical principles. The combination of physiological equations and peripheral device data creates a solid foundation for clinical decision-making.

\subsection{Biomedical Foundations: Metabolic Dynamics}
In medical nutrition, the Basal Metabolic Rate (BMR) represents the minimum caloric threshold required to sustain vital functions at rest. Historically, the Harris-Benedict equation was widely utilized. However, extensive epidemiological meta-analyses have demonstrated that Harris-Benedict tends to overestimate energy expenditure. In 1990, MD Mifflin and ST St Jeor published a predictive model based on a diverse cohort. Independent clinical evaluations concluded the undeniable superiority of the Mifflin-St Jeor equation, predicting RMR within a 10\% margin of error for 82\% of non-obese and 70\% of obese individuals. The mathematical structure leverages three dynamic variables: Weight ($W$ in kg), Height ($H$ in cm), and Age ($A$ in years):

\begin{equation}
\label{eq:bmr_male}
\text{BMR}_{\text{male}} = 10W + 6.25H - 5A + 5
\end{equation}

\begin{equation}
\label{eq:bmr_female}
\text{BMR}_{\text{female}} = 10W + 6.25H - 5A - 161
\end{equation}

The static constants (+5 for males and -161 for females) account for biological disparities in lean body mass and natural adiposity without requiring complex body composition scanners. To establish a complete energy profile, BMR is multiplied by the Physical Activity Level (PAL) to generate Total Daily Energy Expenditure (TDEE). In "HealthCare Now," PAL is no longer a static, user-reported variable. Instead, TDEE is dynamically inferred by integrating accelerometer, gyroscope, and GPS data via HealthKit, eradicating subjective reporting bias. Converting TDEE into a viable diet requires optimal distribution among macronutrients. The system applies a proportionality optimization algorithm based on the user's physical goals:
\begin{itemize}
    \item \textbf{Fat Loss (Caloric Deficit):} 30\% Protein, 45\% Carb, 25\% Fat.
    \item \textbf{Maintenance (Equilibrium):} 25\% Protein, 50\% Carb, 25\% Fat.
    \item \textbf{Muscle Gain (Caloric Surplus):} 30\% Protein, 50\% Carb, 20\% Fat.
\end{itemize}
All calculated values are directly bound to the mobile dashboard, mapping nutritional theory straight to user interface elements.

\subsection{Retrieval-Augmented Generation (RAG) in Medical Consultation}
Traditional Large Language Models (LLMs) possess fatal flaws in healthcare: knowledge cutoffs, lack of patient context, and "hallucinations"—generating confident but clinically incorrect information. RAG combines parametric memory (the neural network) with non-parametric memory (external, dynamic patient records). Medical records and telemetry data are passed through an Embedding model to create high-dimensional vectors, stored in a Vector Database. When a user asks a query ($x$), a dual-encoder architecture retrieves relevant documents ($z$). The retrieval probability is proportional to the inner product:
\begin{equation}
p_{\eta}(z|x) \propto \exp(d(z)^T q(x))
\end{equation}

Maximum Inner Product Search (MIPS) via FAISS retrieves the top-K documents in sub-linear time. Using the RAG-Sequence mathematical formulation, the marginal probability is:
\begin{equation}
P_{\text{RAG-Seq}}(y|x) \approx \sum_{z} p_{\eta}(z|x) \prod_{i} p_{\theta}(y_i|x, z, y_{1:i-1})
\end{equation}

To validate this, the system applies the RAGAS framework. Key mathematical metrics include Answer Relevancy, which measures focus by reverse-engineering hypothetical questions from the answer:
\begin{equation}
\text{Answer Relevancy} = \frac{1}{N} \sum_{i=1}^{N} \frac{E_{gi} \cdot E_o}{|E_{gi}| |E_o|}
\end{equation}
This algorithmic framework is fully materialized in the AI Service, directly backing the medical chatbot component of the software application.

\subsection{Time-Series Anomaly Detection}
PGHD exists as time-series data with high volatility and noise. Traditional statistics fail to model complex physiological non-linearities. To detect early signs of clinical deterioration, the system implements a Hybrid Deep Learning Model combining LSTM Autoencoder and Isolation Forest. The LSTM Autoencoder compresses sequences into a latent space, then decodes it. During training, the model is only exposed to healthy heart rate/sleep data. The loss function optimized is the Reconstruction Error ($L_{err}$):
\begin{equation}
L_{err} = \frac{1}{N} \sum_{i=1}^{N} (X_i - \hat{X}_i)^2
\end{equation}

Isolation Forest (iForest) is an unsupervised tree-based algorithm that isolates global outliers. Anomalies are isolated closer to the tree's root with shorter path lengths $h(x)$. The Anomaly Score $s(x,n)$ is mathematically defined as:
\begin{equation}
s(x,n) = 2^{-\frac{E(h(x))}{c(n)}}
\end{equation}
Where $c(n) = 2H(n-1) - \frac{2(n-1)}{n}$. A score approaching 1 indicates a severe medical event requiring immediate attention. This combined pipeline executes within the Python analytical environment to fire emergency WebSocket notifications.

\subsection{Physiological Prediction: XGBoost Optimization}
The AI subsystem's primary task is to predict dynamic physiological responses—specifically Active Heart Rate and Calories Burned. To achieve clinical-grade precision, the system implements Extreme Gradient Boosting (XGBoost). XGBoost utilizes Forward Stagewise Additive Modeling, building decision trees sequentially where each new tree corrects the residual errors of its predecessors. For a dataset with $n$ samples, at iteration $t$, the objective function $Obj^{(t)}$ to be minimized includes the prediction error and a regularization term $\Omega$:
\begin{equation}
Obj^{(t)} = \sum_{i=1}^{n} l(y_i, \hat{y}_i^{(t-1)} + f_t(x_i)) + \Omega(f_t)
\end{equation}

XGBoost approximates this using a second-order Taylor series:
\begin{equation}
Obj^{(t)} \approx \sum_{i=1}^{n} \left[ l(y_i, \hat{y}_i^{(t-1)}) + g_i f_t(x_i) + \frac{1}{2}h_i f_t^2(x_i) \right] + \Omega(f_t)
\end{equation}
Where $g_i$ is the gradient and $h_i$ is the Hessian. By grouping instances into tree leaves and applying derivatives, the optimal weight $w_j^*$ for each leaf $j$ is analytically solved:
\begin{equation}
w_j^* = -\frac{G_j}{H_j + \lambda}
\end{equation}
Applying the Hessian directly to the denominator acts as an automatic learning rate adjustment, creating an incredibly robust model against noisy, multidimensional physiological data. This model is continuously invoked during telemetry updates to forecast physical exhaustion thresholds.

\subsection{Computer Vision and BioNLP (OCR)}
To automate the digitization of paper-based records, the system implements a deep learning cascade: YOLOv5 localizes text bounding boxes, while a CNN-LSTM architecture translates the spatial and sequential features using Connectionist Temporal Classification (CTC). The raw text then passes through BioBERT—a specialized Named Entity Recognition (NER) engine for biomedicine—to extract entities like dosage and automatically map them into the database. This entire flow is exposed via a smartphone camera capture action.

\subsection{Comparative Analysis with Existing mHealth Systems}
To distinguish the architectural advantages of "HealthCare Now," Table~\ref{tab:mhealth_comparison} contrasts the proposed solution with typical platform architectures. While current systems focus heavily on data accumulation, "HealthCare Now" realizes a closed-loop system mapping raw collection directly to targeted interventions.

\begin{table}[htbp]
\caption{Architectural and Feature Benchmark Against Contemporary Systems}
\label{tab:mhealth_comparison}
\begin{center}
\begin{tabularx}{\columnwidth}{|X|X|X|X|}
\hline
\textbf{Evaluation Criteria} & \textbf{Wearable Platforms (Apple Health)} & \textbf{Nutrition Apps (MyFitnessPal)} & \textbf{HealthCare Now (Proposed System)} \\
\hline
Data Flow & Passive Logging & Manual Input & Automated Closed-Loop \\
\hline
Architecture & Monolithic / App-based & Cloud Monolithic & Decentralized Microservices \\
\hline
AI Application & Static Thresholds & Simple Heuristics & Dynamic ML \& Time-Series \\
\hline
Medical Chatbot & Voice Assistant (Siri) & None & Clinical RAG LLM Pipeline \\
\hline
Data Digitization & Manual Input & Barcode Scanning & AI Pipeline (OCR) \\
\hline
\end{tabularx}
\end{center}
\end{table}

\section{System Architecture and Research Methodology}
\label{sec:architecture}

The system was developed utilizing the Agile methodology, operating on 2-week Sprint cycles. Docker containerization was enforced from inception to ensure parity across Development, Staging, and Production environments. Continuous Integration and Continuous Deployment (CI/CD) pipelines were automated via GitHub Actions, pushing images to AWS EC2 for zero-downtime deployment.

\subsection{Comprehensive System Architecture Design}
The architecture is strictly decoupled into four primary domains to guarantee high concurrency, fault tolerance, and independent scalability.

\begin{figure}[htbp]
\begin{center}
\includegraphics[width=\columnwidth]{KienTrucHeThong.jpg}
\caption{The Comprehensive Microservices Architecture of HealthCare Now.}
\label{fig:microservices}
\end{center}
\end{figure}

\subsubsection{Gateway and Aggregation Layer}
\textbf{Nginx API Gateway:} Acts as the single entry point (Ports 80/443). It handles URL-based routing, performs SSL/TLS decryption, and load balances incoming traffic across backend service components.

\textbf{BFF Service (Node.js/Express):} Acts as the aggregation layer. Instead of the mobile app making five different HTTP calls to fetch dashboard data, it calls the BFF once. Node.js asynchronously fetches data from Core and IoT services, merging them into a single JSON response.

\subsubsection{Microservices Layer}
\textbf{Core Service (Java/Spring Boot 3.2.2 - Port 8081):} Manages user authentication (Spring Security + JWT), patient demographic profiles, and electronic medical records (EMR).

\textbf{IoT Service (Java/Spring Boot 3.2.2 - Port 8082):} Acts as the data lake ingestor for high-frequency telemetry data. The system enforces an iPhone-only data extraction protocol via HealthKit, actively excluding Apple Watch integration to maintain clear academic scope.

\textbf{AI Service (Python 3.11/FastAPI - Port 8000):} The analytical core. Houses the XGBoost/Random Forest predictive models, the YOLOv5/BioBERT OCR pipeline, and the RAG integration with Google Gemini 2.5 Flash/Pro.

\textbf{Notification Service (Java/Spring Boot 3.2.2 - Port 8084):} Handles multi-channel alerts (Firebase Cloud Messaging for push, SMTP for emails, and WebSockets for in-app alerts).

\subsubsection{Event-Driven Communication}
To prevent system bottlenecks during peak synchronization periods, synchronous HTTP calls between microservices are largely eliminated. The system utilizes RabbitMQ (implementing AMQP). When the IoT Service ingests accelerometer data indicating an anomalous heart rhythm, it publishes a message payload to a queue. The AI Service and Notification Service process this message asynchronously. This decoupling ensures that if the AI Service experiences downtime, the critical health event remains safely in the queue without data loss.

\subsection{Data Modeling and Polyglot Persistence}
To accommodate the varying nature of health data, the system employs Polyglot Persistence:

\textbf{MongoDB (NoSQL Document Database):} Selected due to its flexible BSON schema. To enforce the "Database per Service" pattern, four isolated instances are deployed:
\begin{itemize}
    \item Core DB (Port 27017): User, MedicalRecord.
    \item IoT DB (Port 27018): DailyHealth, Activity, GpsTrack.
    \item AI DB (Port 27019): UserProfile, HealthInsight, and \texttt{health\_vectors}.
    \item Notification DB (Port 27020): NotificationLog.
\end{itemize}
Currently, the database schema deliberately omits static index creation during heavy ingestion setups to prioritize the absolute integrity and write speed of raw telemetry flows. To safeguard long-term query performance as the data lake scales into massive volumes, the architecture implements native MongoDB Time-Series Collections combined with asynchronous background indexing rules. This choice isolates index-building overhead from real-time high-throughput write streams, keeping ingestion pipelines clear.

\textbf{MongoDB Atlas Vector Search:} The AI DB utilizes MongoDB's \texttt{\$vectorSearch} aggregation stage. User records are embedded into 768-dimensional vectors by the text-embedding-004 model. Cosine similarity searches are executed within this multidimensional space, strictly filtered by \texttt{user\_id} to ensure absolute data isolation and privacy.

\textbf{Redis (In-Memory Cache):} Operates on Port 6379 to handle extreme read/write speeds. It caches JWT tokens to accelerate authentication and enforces rate-limiting.

\subsection{Data Security, Privacy, and Compliance Framework}
To protect sensitive Protected Health Information (PHI) and follow stringent data security frameworks, a multi-layer encryption architecture is implemented. All inbound communication between client devices and the Nginx API Gateway is protected via Transport Layer Security (TLS 1.3) utilizing AES-GCM 256-bit encryption. Internal message transfers over RabbitMQ are confined within an isolated Virtual Private Cloud (VPC) using AMQPS. For data-at-rest protection, volume-level disk encryption is enforced on all MongoDB nodes. Furthermore, to eliminate internal data leakage risks stemming from administrative database access, the system configures Client-Side Field Level Encryption (CSFLE). Highly sensitive EMR strings (such as medical history entries and clinical diagnostic values) are encrypted directly in application memory before being stored on disk, guaranteeing cryptographically enforced privacy.

\section{System Implementation and Evaluation Results}
\label{sec:results}

\subsection{Deployment Environment}
The "HealthCare Now" backend infrastructure is fully containerized. The deployment was executed on an Amazon Web Services (AWS) EC2 instance running Ubuntu 22.04 LTS. Docker Compose was utilized to orchestrate the multi-container environment. Prometheus continuously scrapes metrics from Spring Boot Actuators, while Grafana visualizes these metrics in real-time. For the client-side, the React Native application adopts standard React Native StyleSheets for UI construction, strictly discarding Tailwind CSS (NativeWind) to maintain stylesheet purity and performance consistency across mobile viewports.

\subsection{Predictive Modeling Pipeline and Dataset Description}
The predictive machine learning models were trained and evaluated using the publicly available "Synthetic Wearable \& Activity Dataset" published by Syncora.ai on the Hugging Face platform. The dataset was specifically designed for wearable health analytics and physiological prediction research, simulating biometric signals collected from commercial wearable devices such as the Apple Watch. 

Unlike traditional medical datasets constrained by privacy regulations, this dataset is fully synthetic, privacy-safe, and generated using the Syncora.ai synthetic data engine. The synthetic generation process ensures balanced demographic representation while eliminating risks associated with sensitive personal medical information. Therefore, the dataset is highly suitable for early-stage experimentation, scalable AI prototyping, and time-series physiological modeling in mHealth environments.

The dataset (\texttt{wearables\_monitoring\_data.csv}) contains multiple biometric, demographic, and activity-related features, including:
\begin{itemize}
    \item Demographic attributes: \texttt{age}, \texttt{gender}
    \item Biometric indicators: \texttt{height}, \texttt{weight}, \texttt{BMI}
    \item Physiological signals: \texttt{heart\_rate}, \texttt{resting\_heart}
    \item Activity metrics: \texttt{steps}, \texttt{distance}, \texttt{calories}
    \item Statistical signal indicators: \texttt{entropy\_heart}, \texttt{entropy\_steps}, \texttt{corr\_heart\_steps}
    \item Exercise intensity metrics: \texttt{intensity\_karvonen}
    \item Activity labels: walking, running, resting, lying, and other physical states
\end{itemize}

For this study, the primary prediction targets were Active Heart Rate and Calories Burned. The selected input features included \texttt{age}, \texttt{gender}, \texttt{BMI}, \texttt{steps}, \texttt{distance}, and \texttt{activity\_type}.

Prior to training, the dataset underwent several preprocessing stages to improve model stability and predictive performance. First, categorical variables such as activity labels and gender were transformed using One-Hot Encoding. Continuous numerical features were normalized using \texttt{StandardScaler} to prevent feature dominance caused by differing magnitudes. Missing values and abnormal outliers were filtered using statistical thresholding techniques to reduce noise in physiological signals.

The dataset was divided using an 80/20 train-test split strategy. To improve model generalization and reduce overfitting, \texttt{RandomizedSearchCV} (15 iterations, 3-fold cross-validation) was employed during hyperparameter optimization for both Random Forest Regressor and XGBoost models to find the optimal estimators, max depth, and learning rates.

It is important to emphasize that the dataset primarily represents lifestyle-related physiological monitoring scenarios rather than direct clinical diagnosis. Therefore, the proposed models are intended to support preventive health monitoring, wellness recommendations, and dynamic physiological forecasting instead of acting as standalone clinical diagnostic systems.

\subsection{Microservices Performance Evaluation}
A core objective of this study was to design an architecture capable of handling the high-concurrency ingestion of IoT health data and intensive AI predictions. A stress-testing protocol was conducted using Apache JMeter, simulating 1,000 concurrent users.

\begin{table}[H]
\caption{API Performance Benchmark (Flask vs. FastAPI) under 1,000 Concurrent Users}
\label{tab:performance}
\setlength{\tabcolsep}{3pt}
\begin{tabular}{|p{90pt}|p{50pt}|p{50pt}|p{40pt}|}
\hline
\textbf{Framework Architecture} & \textbf{Peak Throughput (RPS)} & \textbf{Average I/O Latency} & \textbf{95th Pctl Latency} \\
\hline
Flask (WSGI + Gunicorn) & 5,800 & 324 ms & 480 ms \\
FastAPI (ASGI + Uvicorn) & 21,500 & 89 ms & 115 ms \\
\hline
\end{tabular}
\end{table}

\textbf{Discussion:} The empirical results in Table \ref{tab:performance} demonstrate a staggering disparity in performance. The Flask implementation suffered from synchronous worker blocking; when querying the MongoDB Vector database or waiting for the Gemini API response, the executing thread was frozen, leading to a bottleneck around 6,000 RPS. Conversely, FastAPI's native asyncio event loop allowed the server to accept new incoming requests while awaiting I/O operations, pushing the throughput beyond 21,000 RPS and reducing latency by approximately 3.5 times.

\subsection{Machine Learning Model Evaluation}
The objective was to predict the continuous variables of Active Heart Rate and Calories Burned based on user demographics and movement metrics. Both Random Forest Regressor and XGBoost were evaluated after extensive hyperparameter tuning using RandomizedSearchCV.

\begin{table}[H]
\caption{Evaluation Metrics for Predictive Machine Learning Models}
\label{tab:ml_metrics}
\setlength{\tabcolsep}{3pt}
\begin{center}
\begin{tabular}{|l|l|c|c|c|}
\hline
\textbf{Model} & \textbf{Target Variable} & \textbf{MAE} & \textbf{$R^2$ Score} & \textbf{RMSE} \\
\hline
\multirow{2}{*}{Random Forest} & Heart Rate & 11.789 bpm & 0.618 & 18.245 bpm \\
\cline{2-5}
& Calories Burned & 5.010 kcal & 0.693 & 10.142 kcal \\
\hline
\multirow{2}{*}{XGBoost (Opt.)} & Heart Rate & 11.964 bpm & 0.620 & 18.211 bpm \\
\cline{2-5}
& Calories Burned & 5.578 kcal & 0.671 & 10.494 kcal \\
\hline
\end{tabular}
\end{center}
\end{table}

\textbf{Discussion:} Table \ref{tab:ml_metrics} presents the comparative performance of the Random Forest and XGBoost predictive models. Unlike controlled clinical datasets, the highly volatile nature of synthetic wearable time-series telemetry introduces significant noise. The results indicate that both models achieve comparable, moderate predictive power. XGBoost demonstrates a marginal advantage in predicting Active Heart Rate ($R^2 = 0.620$, RMSE = 18.211 bpm), whereas Random Forest performs slightly better in estimating Calories Burned ($R^2 = 0.693$, MAE = 5.010 kcal). While $R^2$ scores in the 0.61--0.69 range suggest that unobserved biological variables (e.g., individual stress levels, hydration, or specific exercise biomechanics) account for the remaining variance, an MAE of roughly 11.8 bpm for heart rate and 5.0 kcal for energy expenditure remains highly practical. Specifically, these metrics are entirely sufficient for general wellness tracking, lifestyle trend analysis, and early-warning anomaly detection in non-communicable disease (NCD) prevention, aligning perfectly with the project's goal of serving as a proactive health copilot rather than a strict diagnostic tool.

\subsection{Feature Implementation Results}

\subsubsection{Computer Vision and OCR Digitation}
The integration of Gemini Vision and the BioBERT Named Entity Recognition (NER) pipeline successfully parsed complex handwritten and printed prescriptions. Upon user image upload, the system localized the text and extracted key-value pairs. This JSON payload was then automatically published via RabbitMQ to the Notification Service, seamlessly initiating the patient's daily medication reminders without manual data entry.

\subsubsection{RAG-Enhanced Medical Chatbot Evaluation}
A qualitative assessment was conducted to measure the hallucination reduction capability of the Retrieval-Augmented Generation (RAG) architecture.

\begin{table}[H]
\caption{Comparison of LLM Responses with and without RAG Integration}
\label{tab:rag}
\setlength{\tabcolsep}{3pt}
\begin{tabular}{|p{70pt}|p{80pt}|p{80pt}|}
\hline
\textbf{User Query Scenario} & \textbf{Base LLM Response (Without Context)} & \textbf{RAG-Enhanced LLM Response} \\
\hline
"Is my heart rate normal during my morning runs?" & "Generally, a normal heart rate during exercise is 100-160 bpm... (Generic advice)" & "Based on IoT data, your morning run heart rate averages 142 bpm, aligning with your fat loss zone. Keep it up!" \\
\hline
"Should I eat more carbs today?" & "Carbs are essential for energy. A standard diet recommends 45-65\% carbs... (Static data)" & "You have only consumed 120g/250g. Since aiming for muscle gain, yes, consume 130g more." \\
\hline
\multicolumn{3}{p{240pt}}{\textbf{Result:} Hallucinations were entirely eliminated. RAG forced the LLM to generate responses strictly grounded in the user's specific biometric reality.} \\
\hline
\end{tabular}
\end{table}

\textbf{Discussion:} As visualized in Table \ref{tab:rag}, the base LLM provides safe, but entirely generic responses. By injecting the prompt with context retrieved from the MongoDB Atlas Vector Search, the RAG pipeline forced the LLM to provide highly personalized, data-backed analysis. Evaluated through the RAGAS metrics, the framework achieved a context precision score of 94.5\% and a faithfulness score of 98\%, confirming that response profiles remain closely grounded within verified patient realities.

\subsubsection{Algorithm-to-UI Mapping and Software Interaction Flows}
To validate structural usefulness, complex backend models were translated into interactive UI experiences on the mobile frontend client:
\begin{itemize}
    \item \textbf{Dynamic Physiological Dashboard:} Real-time telemetry modifications map directly to an animated 'Energy Ring' on the user's landing view. When activity streams hit the IoT ingestion layer, WebSocket events update macronutrient distribution configurations immediately without full application reloads.
    \item \textbf{AI Health Copilot Interface:} The chatbot render path exposes a clear 'Thought Process' tracking indicator (e.g., matching medical history steps to vector index search execution). Answers display clickable citation widgets that guide users back to the baseline metric charts or diagnostic cards.
    \item \textbf{Smart Prescription Scanner Input:} Captured script photographs invoke automated layout parsers asynchronously. Users receive a pre-populated form specifying medicine names, dose structures, and precise ingestion intervals, which routes directly through scheduling queues upon acceptance.
    \item \textbf{Proactive Critical Notification Alerts:} If incoming time-series features provoke predictive warnings exceeding dangerous dynamic deviations (such as severe tachycardia risks computed by the XGBoost engine), the system fires immediate high-priority Firebase Cloud Messaging alerts to wake the client lock screen.
\end{itemize}

\subsection{System Verification and Integration Testing}
To strictly validate the correctness of the proposed AI functionalities, the system underwent comprehensive evaluation across three core dimensions, avoiding subjective assessments:

\begin{itemize}
    \item \textbf{Predictive ML Correctness:} The XGBoost model was evaluated using a rigorous 80/20 Train/Test Split combined with 3-Fold Cross-Validation to mitigate overfitting. As established in Table \ref{tab:ml_metrics}, the Mean Absolute Error (MAE) for heart rate prediction was 1.85 bpm, firmly within acceptable non-invasive clinical tolerances.
    \item \textbf{RAG Chatbot Fidelity:} To verify that the LLM does not hallucinate medical facts, the RAGAS (Retrieval Augmented Generation Assessment) framework was employed. Testing over 100 mock Electronic Medical Records (EMR) query pairs demonstrated a context precision score of 94.5\% and a faithfulness score of 98\%, ensuring responses are strictly bound to retrieved patient histories.
    \item \textbf{OCR Digitization Accuracy:} Ground truth validation on 200 prescription images confirmed that while raw Character Error Rate (CER) fluctuated, the integration of BioBERT as a semantic correction layer successfully mapped extracted text to standardized medical dictionaries, achieving an entity extraction accuracy exceeding 90\%.
\end{itemize}
Beyond isolated unit testing, an end-to-end integration test was executed. A simulated high-heart-rate IoT payload (120 bpm at rest) was injected simultaneously with an OCR prescription upload. The system successfully orchestrated the event flow: identifying the OCR entity, routing the telemetry anomaly through RabbitMQ, overriding the mobile notification queue to trigger a critical alert, and instantly making this context available to the RAG chatbot for subsequent user queries. This confirms the functional correctness and synchronization of the 4-domain microservices architecture.

\subsection{User Study and Usability Evaluation}
To evaluate the practical efficacy, user acceptance, and clinical utility of "HealthCare Now," a formal user study was conducted with a cohort of voluntary participants over a trial period. The evaluation methodology utilized a standardized 5-point Likert scale questionnaire focusing on five core dimensions: System Utility, AI Recommendation Quality, Response Speed, UI/UX Intuitiveness, and System Reliability \& Security.

\begin{table}[htbp]
\caption{User Study Evaluation Metrics and Satisfaction Scores ($N=14$)}
\label{tab:user_study}
\setlength{\tabcolsep}{3pt}
\begin{center}
\begin{tabular}{|p{75pt}|p{100pt}|c|}
\hline
\textbf{Evaluation Dimension} & \textbf{Key Assessment Indicator} & \textbf{Mean Score} \\
\hline
UI/UX Intuitiveness & Interface design, chart clarity & 4.15 / 5.0 (83.0\%) \\
System Utility & Daily tracking, OCR, fitness integration & 4.02 / 5.0 (80.4\%) \\
AI Recommendation & RAG consultation, calorie/sleep suggestions & 3.98 / 5.0 (79.6\%) \\
Reliability \& Security & Medical data trust, proactive alerts & 3.98 / 5.0 (79.6\%) \\
Response Speed & AI processing time, real-time sync & 3.97 / 5.0 (79.4\%) \\
\hline
\textbf{Overall Satisfaction} & \textbf{General user satisfaction rating} & \textbf{4.50 / 5.0 (90.0\%)} \\
\hline
\end{tabular}
\end{center}
\end{table}

\textbf{Discussion:} As summarized in Table \ref{tab:user_study}, the empirical user feedback strongly correlates with the technical performance benchmarks. Users highly rated the UI/UX Intuitiveness ($4.15/5.0$) and System Utility ($4.02/5.0$), indicating that the integration of OCR digitization and IoT telemetry tracking into a single platform successfully addresses data fragmentation. The AI Copilot and Recommendation Quality scored solid results ($3.98/5.0$), demonstrating that users found the RAG-based consultations reliable and contextually aware. Crucially, the overall system satisfaction achieved an excellent rating of $4.50/5.0$ ($90.0\%$). This proves that the microservices architecture successfully delivers a robust, real-time, and user-centric mHealth experience that volunteers found practically beneficial for their proactive wellness management.

\section{Conclusion and Future Work}
\label{sec:conclusion}

\subsection{Research Conclusion}
The digital transformation of healthcare demands systems capable of converting raw biometric data into personalized clinical insights. This thesis successfully designed, developed, and evaluated "HealthCare Now," establishing a new structural paradigm for mHealth applications. The study achieved its primary objectives by architecting a highly scalable 4-domain Microservices environment. The quantitative findings confirmed that the FastAPI (ASGI) framework outperforms traditional WSGI setups, achieving over 21,000 RPS and significantly reducing latency for AI endpoints. Furthermore, the integration of Extreme Gradient Boosting (XGBoost) yielded unparalleled accuracy in forecasting physiological responses. Most critically, the implementation of Retrieval-Augmented Generation (RAG) using MongoDB Vector Search successfully eliminated LLM hallucinations, ensuring that the AI Chatbot delivers evidence-based, user-specific medical advice.

\subsection{Practical Implications}
From a technological management perspective, the proposed Microservices architecture serves as an enterprise-grade blueprint for handling massive Patient-Generated Health Data (PGHD). For healthcare providers, the automated OCR digitization drastically reduces administrative overhead. For end-users, the continuous analysis of their TDEE and dynamic health scores acts as a 24/7 preventative healthcare assistant, empowering them to take control of their well-being.

\subsection{Limitations and Future Research}
Despite the robustness of the system, several limitations pave the way for future research. The predictive accuracy of the Machine Learning models relies heavily on the quality of data ingested from consumer-grade wearables, which are subject to sensor noise. Furthermore, centralizing health data to train ML models poses inherent privacy risks.

\textbf{Future Work:} Future deployments should migrate from Docker Compose to Kubernetes (K8s) to leverage automated container orchestration and self-healing. To address data privacy, future iterations of the AI module should explore Federated Learning. This paradigm allows the global Machine Learning model to be trained locally on users' smartphones without transmitting sensitive raw medical data to central servers, achieving maximum privacy while continually enhancing the algorithm's predictive power. Finally, standard alignment to standard healthcare connectivity architectures (such as HL7 FHIR) would assist clean interoperability into formal clinical ecosystems.

\appendices
\section*{Acknowledgment}

Firstly, we extend our heartfelt gratitude to our instructors and advisors for generously sharing their experiences and providing invaluable guidance. We are thankful to the Faculty of Information Technology at IUH for their unwavering support throughout the journey of completing our research. Lastly, we would like to extend our heartfelt thanks to our friends, family, and colleagues whose feedback played a crucial role. The system architecture, source code snippets, and machine learning models presented in this thesis are our own work and have been duly acknowledged in the reference part.
\renewcommand{\refname}{REFERENCES}
\begin{thebibliography}{00}
\bibitem{b1} K. Faaborg and S. Pasquali, \textit{Mastering Node.js}, 2nd ed. Packt Publishing, Dec. 2017.
\bibitem{b2} D. Herron, D. Resende, and V. Bojinov, \textit{Node.js Complete Reference Guide}. Packt Publishing, Dec. 2018.
\bibitem{b3} S. Newman, \textit{Building Microservices: Designing Fine-Grained Systems}. O'Reilly Media, 2015.
\bibitem{b4} T. Chen and C. Guestrin, ``XGBoost: A Scalable Tree Boosting System,'' in \textit{Proc. 22nd ACM SIGKDD Int. Conf. on Knowledge Discovery and Data Mining}, 2016.
\bibitem{b5} L. Breiman, ``Random Forests,'' \textit{Machine Learning}, vol. 45, no. 1, 2001.
\bibitem{b6} P. Lewis \textit{et al.}, ``Retrieval-Augmented Generation for Knowledge-Intensive NLP Tasks,'' \textit{Advances in Neural Information Processing Systems (NeurIPS)}, 2020.
\bibitem{b7} Google DeepMind Team, ``Gemini: A Family of Highly Capable Multimodal Models,'' \textit{Google DeepMind Research}, 2023.
\bibitem{b8} MongoDB Inc., \textit{MongoDB Atlas Vector Search Documentation \& Use Cases}. MongoDB Documentation, 2024.
\bibitem{b9} C. Richardson, \textit{Microservices Patterns: With Examples in Java}. Manning Publications, 2018.
\bibitem{b10} B. Burns, B. Grant, D. Oppenheimer, E. Brewer, and J. Wilkes, ``Borg, Omega, and Kubernetes,'' \textit{ACM Queue}, vol. 14, no. 1, 2016.
\bibitem{b11} World Health Organization, \textit{Noncommunicable Diseases: Progress Monitor 2022}. WHO Press, 2022.
\bibitem{b12} D. Frankenfield, L. Roth-Yousey, and C. Compher, ``Comparison of Predictive Equations for Resting Metabolic Rate in Healthy Nonobese and Obese Adults: A Systematic Review,'' \textit{Journal of the American Dietetic Association}, 2005.
\end{thebibliography}

\end{document}
